%% Owner: Víctor Martín Castilla
\section{Architecture}
\label{sec:architecture}

The statement fixes the services and groups them into stacks
\parencite{llusa2026project}. Each stack is one Compose file deployed as
one Swarm stack, so it can be developed, deployed and rolled back on its
own \parencite{llusa2026compose,llusa2026orchestration}.

\begin{description}
  \item[basic] Shared services. An nginx front proxy receives all public
    traffic and forwards it to the other stacks. Certbot obtains and
    renews Let's Encrypt certificates for it. A private registry stores
    the images the team builds, and an NFS export provides the disk
    shared between nodes.
  \item[app] The company's core application, a Python web service, with
    its PostgreSQL database, a Redis cache, an Apache server for static
    files, an SMTP relay towards an external mail provider, an automatic
    feeder that captures ratings and posts them to the application, and a
    cron job that dumps the database. \todo{describe what the core app
    does once it exists.}
  \item[atenea] The Moodle learning management system with its MariaDB
    database, the Moodle cron and a backup job.
  \item[stats] Grafana, provisioned from files following the
    observability-as-code approach, with Prometheus as its metrics store,
    node exporter and cAdvisor on every node, and a Grafana plugin.
    \todo{name the plugin and justify the choice.}
\end{description}

\Cref{tab:services} lists every service with its number of replicas and
where Swarm may place it. The statement requires every service to scale
to several containers, except cron jobs and database engines that do not
allow it \parencite{llusa2026project}. Every single-replica row falls in
one of those two groups. Prometheus counts as a database because it keeps
its time series on local disk. Redis counts as one because two
independent instances behind one virtual IP would hold two different
caches.

\begin{table}[htbp]
  \centering
  \caption{Services, replicas and placement. ``Any'' means Swarm chooses
    the node; \texttt{role=\ldots} pins the service to the node carrying
    that label (\cref{sec:swarm}).}
  \label{tab:services}
  \small
  \begin{tabular}{lllll}
    \toprule
    Stack  & Service          & Replicas  & Placement                    & Reason \\
    \midrule
    basic  & nginx front      & 2 or more & any                          & stateless \\
           & certbot          & 1         & \texttt{role=entry}          & cron job \\
           & registry         & 2         & any, storage on NFS          & shared storage \\
           & off-site copy    & 1         & \texttt{role=entry}          & cron job \\
    \midrule
    app    & core app         & 2 or more & any                          & stateless \\
           & PostgreSQL       & 1         & \texttt{role=db}             & DBMS \\
           & Redis            & 1         & any                          & DBMS, cache only \\
           & static Apache    & 2 or more & any                          & stateless \\
           & SMTP relay       & 2         & any                          & stateless \\
           & auto feeder      & 2         & any                          & stateless \\
           & database dump    & 1         & any                          & cron job \\
    \midrule
    atenea & Moodle           & 2 or more & any, files on NFS            & shared storage \\
           & MariaDB          & 1         & \texttt{role=db}             & DBMS \\
           & Moodle cron      & 1         & any                          & cron job \\
           & Moodle backup    & 1         & any                          & cron job \\
    \midrule
    stats  & Grafana          & 2         & any                          & shared database \\
           & Grafana database & 1         & \texttt{role=monitoring}     & DBMS \\
           & Prometheus       & 1         & \texttt{role=monitoring}     & time-series DBMS \\
           & node exporter    & global    & every node                   & one per host \\
           & cAdvisor         & global    & every node                   & one per host \\
    \bottomrule
  \end{tabular}
\end{table}

One component is not a container. The NFS server runs on the host
system of one node, because a kernel NFS server needs privileged mode,
which Swarm services do not support.
